\documentclass[UTF8]{ctexart}
\usepackage{amsmath}
\usepackage{pdfpages}
\usepackage{tocloft}
\usepackage{amsthm,amssymb}
\usepackage{booktabs}
\usepackage{graphicx}
\usepackage{siunitx}
\usepackage[hidelinks]{hyperref}
\usepackage{geometry}
\usepackage{float}
\usepackage{multirow}
\usepackage{indentfirst}
\usepackage{setspace}
\usepackage{fancyhdr}
\usepackage{diagbox}
\usepackage{subcaption}
\usepackage{makecell} 
\usepackage{listings}
\usepackage{xcolor}
\usepackage{tikz}
\usepackage{circuitikz}
\usepackage{longtable}

\renewcommand{\cftsecleader}{\cftdotfill{\cftdotsep}}
\newcommand\question[2]{\noindent\fbox{\parbox[t]{\linewidth}{\hspace{0pt}{\textbf{#1\quad}#2}}}}
\renewcommand\appendix{\par
	\setcounter{section}{0}
	\setcounter{subsection}{0}
	\gdef\thesection{附加说明}}
\geometry{a4paper,left=2cm,right=2cm,top=2cm,bottom=2cm} 
\linespread{1.5}
\CTEXsetup[format={\Large \bfseries}]{section}
\title{\vspace{-2em}\heiti 偏振编码的BB84协议量子密钥分发\\实验报告}
% \author{\kaishu C2组\\\kaishu 董皓彧（无37，2023010659）\\\kaishu 刘佳起（无37，2023010668）\\ \kaishu 周王昊（无35，2023010608）}
\date{}

\begin{document}
\maketitle

\begin{abstract}

量子密钥分发（QKD）旨在让窃听者获取尽量少信息的情况下实现密钥的安全产生和分配。本次实验中我们搭建了一套小型偏振编码BB84协议QKD通信系统，通过偏振片与波片产生不同偏振态的光用于编码信息，经过衰减与分光后信号被单光子探测器接收，在输出端进行光子计数与数据处理，再根据BB84协议步骤通过经典信道交换信息产生初始密钥。实验测得输入端与输出端得到的初始密钥能够达到90\% 以上的匹配率，使用LDPC+Cascade数据后处理能够生成安全密钥。

\end{abstract}

\section{引言}

以量子密钥分发（QKD）为核心的量子保密通信技术已发展到与工程技术结合向实用化发展的阶段。QKD 将单光子序列的某一量子态（如偏振态）编码为密钥，利用海森堡不确定性原理和量子态不可克隆原理实现密钥的安全产生和分配，在此过程中窃听者无法获知完整准确的密钥信息，且窃听行为可以被检测到。1984年，Bennett等人提出了第一个QKD协议，称为BB84协议。对于自由空间量子信道，光的偏振态便于产生、传送和探测，因此偏振编码的BB84协议QKD是应用最为广泛的QKD协议之一。

为了通过实验加深对光量子信息技术和QKD原理的理解，更好地掌握光学和电学系统和相关设计设计与搭建方法，我们使用实验室中的光学仪器搭建了一套小型偏振编码BB84协议QKD通信系统。经过光路调试、发射端驱动电路搭建、衰减与遮光、输出端解码程序编写及系统联调等实验步骤，最终输入端与输出端得到的初始密钥能够达到90\% 以上的匹配率，QBER低于QKD安全性理论阈值，能够通过数据后处理生成安全密钥，证明实验中搭建的QKD系统能够正常实现功能。

\section{基本原理}

\subsection{偏振编码原理与BB84协议通信过程}

偏振编码，即使用光的偏振态编码信息。BB84协议QKD中采用直角基和对角基两组偏振基，直角基中水平偏振代表0，垂直偏振代表1；对角基中$-\ang{45}$代表0，$+\ang{45}$代表1。根据量子力学原理，每个光子在被正确的偏振基接收时以100\%概率被解码为正确信息，而被错误的偏振基接收时各有50\%的概率被解码为0或1，如下图所示。

\begin{figure}[H]
	\centering
	\includegraphics[width=14cm]{graph1.png}
	\caption{偏振编码原理图}
\end{figure}

BB84协议的目的是在发射端和接收端生成相同的一串比特序列作为密钥且使第三方监听者能够获得尽可能少的信息，其通信过程需要两条信道：量子信道（即光路）与经典信道（如计算机网络通信）。发射端Alice首先生成一列随机比特序列作为信息，并为每个信息位随机选择一种偏振基，由信息位和偏振基可确定每个发射单光子波包的偏振方向，通过量子信道将相应偏振方向的单光子波包序列发射给接收端Bob。Bob接收到信息后对每一个单光子波包随机等概率地采用两种测量基中的一种进行测量，得到测量结果序列。由于接收基是随机选择的，测量结果会有一定错误，通信双方需要通过经典信道进行密钥协商。Bob将每个测到光子时采用的测量基序列告知Alice，Alice通过比对测量基序列和单光子波包发射时选取的发射基序列，即可确定两者一致的位置，并通过经典信道告知Bob。两者将发射基与接收基能够匹配的位置对应的比特信息保留下来，形成的比特序列即为原始密钥。

采用两套偏振基、两条信道的通信办法是为了保证通信的安全性。Bennett在关于BB84协议的初始论文中讨论了窃听者Eve采用截取-重发模型进行攻击时对系统误码率的影响，由于Eve并不知道发射基序列，其一定无法获得完整的密钥信息；即使其采用最优的测量基进行测量，即$+\ang{22.5}$方向的偏振基，能够达到85\%的信息准确率，但是根据海森堡不确定性原理和量子态不可克隆原理，Eve无法得知光子原始的偏振方向，只能按照自己测量的结果发射一个相同偏振方向的光子，这样Bob端会检测到25\%的错误概率，能够让其清楚地意识到出现了攻击。因此窃听者能够获得的信息必然是有限的，保证了量子通信的安全性。

\subsection{数据后处理原理}

经过基比对后，Alice 与 Bob 持有的原始密钥序列 $x_a, x_b$ 仍存在一定比例的不一致比特（QBER），需要通过数据后处理进行纠错与隐私放大，使双方最终持有完全一致的安全密钥。本实验中采用低密度奇偶校验码（LDPC）配合多轮披露（类 Cascade）的混合方案：LDPC 负责一次性消除大部分错误，多轮披露用于补救残余错误。

\begin{figure}[H]
	\centering
	\includegraphics[width=14cm]{graph6.png}
	\caption{LDPC+Cascade纠错示意图}
\end{figure}

如图，在 LDPC + Cascade 纠错过程中，发送端与接收端分别将纠错前密钥分组并与校验矩阵 $H$ 相乘，获得校验子 $s_a=H\cdot x_a$ 与 $s_b=H\cdot x_b$，两端通过经典信道交换校验子，比对得到 $s=s_a\oplus s_b$。若本轮解码成功，记错误向量 $e=H^{-1}s$（实际通过置信传播迭代求解），则 $e=x_a\oplus x_b$，Bob 用 $e$ 翻转自身比特即可与 Alice 一致。若本轮未解码成功，则通过经典信道公开 $x_a$、$x_b$ 中若干最不确定的比特位的真值，将这些位作为 \emph{shortened bits}（已知位）并入下一轮的 BP 迭代，直到收敛或达到最大轮数。

\textbf{LDPC 的优势}\quad 校验矩阵 $H$ 稀疏，迭代解码（如归一化最小和 BP 算法）的计算复杂度仅与 $H$ 中非零元个数成线性关系，适合实时纠错。同时通过预先打孔（puncturing，发送方不传输的位，接收方按擦除处理）和缩短（shortening，发送方与接收方约定固定为 0 的位）可以构造任意码率，匹配实际 QBER。

\textbf{Cascade 式多轮披露的作用}\quad 在 QBER 较高、单次 LDPC 解码无法收敛时，公开少数最不确定比特的真值能显著降低剩余错误模式的不确定度，使后续 BP 迭代更容易收敛，从而以较小的额外信息泄漏换取纠错成功率的提升。

\textbf{安全性}\quad 后处理阶段在经典信道上泄漏的信息量主要由每个分组的校验子长度与各轮披露的比特数构成。根据 Shannon 第一定理，纠错任务所需泄漏的信息量存在一个下界（即 $N\cdot h_2(\text{QBER})$，其中 $h_2$ 为二元熵函数）；实际泄漏量越接近这个下界，纠错效率越高。在后续的隐私放大阶段，通过通用哈希函数将纠错后密钥按已泄漏的信息量进行压缩，使得窃听者从经典信道获得的信息无法对最终密钥提供任何有效推断，实现信息论意义上的安全。
\section{实验装置与方法}

\subsection{实验系统总体概述}

\begin{figure}[H]
	\centering
	\includegraphics[width=14cm]{graph2.png}
	\caption{实验系统示意图}
\end{figure}

如图，发射端控制电路控制四个激光器轮流发射激光信号，偏振方向标定为水平和垂直的偏振片用于产生直角基偏振光，半波片用于将直角基偏振光的偏振方向旋转$\ang{45}$变为对角基偏振光。偏振分束镜能够透过水平偏振光并反射垂直偏振光，用于合成或分开不同偏振态的光；而1：1分束镜能够不加选择性地合光与分光，当光被衰减片衰减至单光子量级，在接收端经过1：1分束镜时会等概率随机透射或反射，实现“随机选择测量基”的效果。光子经过与发射光路结构对称的接收光路被单光子探测器接收，由单光子计数器计数并进入后续的数据处理阶段。

实验系统的实拍照片如下：

\begin{figure}[H]
	\centering
	\includegraphics[width=12cm]{graph3.png}
	\caption{实验系统实物图}
\end{figure}

\subsection{实验步骤与方法介绍}

\subsubsection{基础光路搭建与调节}

实验室中的大型激光器偏振方向为垂直方向，而小激光器发光没有特定的偏振方向。旋转偏振片角度直至大激光器发出的光经过偏振片消光，即标定了偏振片方向为水平方向；再将另一个偏振片置于已标定水平方向的偏振片之后，当小激光器发出光束通过两个偏振片后消光，表明第二个偏振片被标定为垂直方向。波片方向的标定需要将一个已标定的偏振片旋转$\ang{45}$角得到$\ang{45}$方向偏振片，旋转波片方向，直至水平或垂直方向偏振光经过波片后能够被$\ang{45}$方向偏振片消光，说明波片起到了将偏振方向旋转$\ang{45}$的效果。

使用上述方法标定各个偏振片与波片方向后开始进行Alice端光路组建，首先使用光屏调节确保各个器件等高，随后微调四个激光器的位置与角度，使四束光从Alice端出射时汇聚在一起。由于小激光器发出的光并非自然光，具有一定的偏振选择性，因此四路光汇聚时的光强有一定区别。以输出时光强最小的一路为基准，我们在其余三路的激光器前各放置了一个偏振片削减光强，使四路光在Alice端出射时的光强基本一致。

由于输入端已经完成光路对齐，Bob端光路调节对等高共轴的要求会低很多，只要确保光路能够从四路出射就行。Bob端波片方向标定可以从输出光功率的角度考虑，实验光路系统的根本目的是接收端使用正确的基测量时偏振方向正确一端能够接收到最大的光功率，而错误一端接收到最小光功率。因此可以在$+\ang{45}$一路发射激光，输出端对称位置探测光功率，旋转波片，当探测光功率最大时，即表明波片正确地旋转了光的偏振方向。

想要实现最单光子的计数，还需要调节接收光电探测器的位置与角度使其与光路耦合。实验时发现输出端耦合的精度要求非常高，即使光电探测器方向出现一点偏差也会对耦合效率产生很大影响，因此需要时常调节。耦合调节的具体方法是使用激光笔从光纤输出端反向向光电探测器射入光，根据光路可逆原理，光纤应该按照原路返回发射端的四个激光器，据此可以对光电探测器的位置与角度进行粗调节。再从输入端射入激光，耦合基本调节完毕时，光纤末端能够输出较强的光功率，因此可以使用光功率计探测输出功率，微调光电探测器的位置与角度使输出功率最大，完成耦合调节。

\subsubsection{输入端控制电路与驱动电路}

\textbf{FPGA 控制逻辑}\quad 发射端使用一块基于 Artix-7 系列芯片（xc7a35tfgg484-2，系统时钟 100 MHz）的 FPGA 完成三项功能：（1）产生频率可调的同步脉冲，输出给单光子计数器（TCSPC）作为时间基准；（2）通过 UART 串口（波特率 250 kbps）从上位机接收 Alice 程序产生的随机字节流；（3）在每个同步上升沿到来时，按照接收到的字节驱动四路激光器中的一路发光。

控制逻辑用 Verilog 实现，由顶层模块 \texttt{top.v} 统一调度四个子模块：\texttt{sync\_gen.v} 实现可调分频的同步信号发生器，默认设置同步频率 50 kHz、脉宽 10 µs；\texttt{uart\_rx.v} 是 8N1 协议的 UART 接收器，含两级同步器规避亚稳态，输出一拍宽度的 \texttt{rx\_valid} 脉冲；\texttt{laser\_ctrl.v} 是激光器驱动控制器，根据 2 bit 的偏振选择信号以独热（one-hot）方式驱动相应通道并维持固定脉宽 10 µs；顶层模块本身包含一个 4 状态主控有限状态机：

\begin{enumerate}
    \item \texttt{S\_WAIT\_RX}：等待 UART 收到 1 个字节；
    \item \texttt{S\_CHECK\_Q}：若该字节为 \texttt{0x71}（即 ASCII \texttt{'q'}）则进入停止态，否则继续；
    \item \texttt{S\_WAIT\_SYNC}：等待 \texttt{sync\_out} 的下一个上升沿，触发 \texttt{laser\_ctrl} 发射一次脉冲后回到 \texttt{S\_WAIT\_RX}；
    \item \texttt{S\_STOPPED}：进入后须按复位键才能重启。
\end{enumerate}

约定每字节的低 2 位编码偏振信息：\texttt{00} 对应 H（0°），\texttt{01} 对应 V（90°），\texttt{10} 对应 +45°，\texttt{11} 对应 -45°，与后文表中的发送态编号保持一致。上位机端 \texttt{Alice/random-gen.py} 一次性产生约 $T\cdot f$ 个随机字节（$T$ 为实验时长，$f$ 为同步频率 50 kHz），通过串口连续发送，最后追加一个 \texttt{'q'} 字节让 FPGA 进入停止态，同时将原始随机序列存盘供后续基比对使用。

\textbf{驱动电路}\quad 由于 FPGA 的输出电流过低，控制激光器输出光强很小，不方便光路的调试，因此我们在 FPGA 之后添加了一个运算放大器开关电路，在控制激光器亮灭的同时使激光器能够输出更高的光强。其原理图与实物图分别如下所示：

\begin{figure}[H]
    \centering
    \begin{circuitikz}[american, thick]
        % ===== 关键节点（绝对坐标）=====
        % 运放置于原点；MOSFET 中心 (3.5, 0)；门极正好在 y=0，与运放输出对齐
        \coordinate (GPIO)  at (-3.5, 0.25);   % GPIO 输入标签 (会平推到 .+ 的 y)
        \coordinate (Vdd)   at (3.7, 2.7);     % Vdd 顶端标签
        \coordinate (mVcc)  at (0, -0.95);     % -Vcc 标签
        \coordinate (Vcc)   at (0, 2.0);       % +Vcc 标签（拉远，留出 LF353 位置）
        \coordinate (GND)   at (3.7, -3.3);    % 接地端
        \coordinate (FB_y)  at (0, -1.6);      % 反馈水平线所在的 y

        % ===== 运放（LF353）——使用 noinv input up 让 + 在上、− 在下 =====
        \draw (0,0) node[op amp, noinv input up] (OA) {};
        % LF353 放在运放正上方、V_cc 引线左侧的空白处（默认 op-amp 高约 2 单位，故 y>1 才在三角形外）
        \node[font=\scriptsize] at (-0.6, 1.35) {LF353};

        % 电源
        \draw (OA.up)   -- (OA.up |- Vcc)  node[above, font=\footnotesize] {$V_{cc}{=}{+}12\,\mathrm{V}$};
        \draw (OA.down) -- (OA.down |- mVcc) node[below, font=\footnotesize] {$-V_{cc}$};

        % GPIO → + 输入（用相对位移 ++(-3, 0) 强制水平）
        \draw (OA.+) -- ++(-3, 0) node[left, font=\footnotesize] {GPIO};

        % ===== MOSFET（2N7000，N-MOS）=====
        \node[nmos] (M) at (3.5, 0) {};
        \node[anchor=west, xshift=8pt, font=\scriptsize] at (M.east) {2N7000};

        % 运放输出 → MOSFET 栅极（两端 y=0，纯水平）
        \draw (OA.out) -- (M.gate);

        % ===== Vdd → D1 → 漏极（路径方向向下，使 anode 在上、cathode 在下，正向偏置）=====
        \draw (M.drain) -- ++(0, 0.4) coordinate (DL);
        \coordinate (DH) at ($(DL) + (0, 1.4)$);
        \draw (DH) to[D, l_=$D_1$] (DL);
        \draw (DH) -- (DH |- Vdd) node[above, font=\footnotesize] {$V_{dd}{=}5/12\,\mathrm{V}$};

        % ===== 源极 → R1 → GND，反馈分叉点 FBJ =====
        \draw (M.source) -- ++(0, -1.1) coordinate (FBJ) node[circ] {};
        \draw (FBJ) to[R, l_=$R_1$] (FBJ |- GND) node[ground] {};

        % ===== 反馈：FBJ → 水平左行 → 垂直上行 → OA.- =====
        \draw (FBJ) -| (OA.-);
    \end{circuitikz}
    \caption{激光器驱动电路原理图（单通道）}
    \label{fig:driver-schematic}
\end{figure}

电路工作原理：FPGA 引脚输出方波 \texttt{IN} 接 LF353 同相输入端，反相输入端从 MOSFET 源极采样电流取样电阻 $R_s$ 上的压降，构成深度负反馈。当 \texttt{IN} 为高时，运放强制 $V_{R_s}\approx V_{\text{IN}}$，恒定电流 $V_{\text{IN}}/R_s$ 流过激光二极管驱动其发光；当 \texttt{IN} 为低时，运放输出钳位在负电源附近，MOSFET 关断，激光器熄灭。该电路本质上是一个由 FPGA 电平开关的恒流源，能在不超过激光器最大电流的前提下提供远高于 FPGA 直接驱动的光强。

\begin{figure}[H]
	\centering
	\includegraphics[width=10cm]{graph4.png}
	\caption{激光器驱动电路实物图（四通道）}
\end{figure}

\subsubsection{衰减与遮光}

为了能够使输出端能够准确地进行解码，需要尽量避免在一个周期内接收到多个光子，一般控制每个周期平均接收到$\eta=0.1\sim0.2$个光子。取$\eta=0.1$进行估算，控制端发射信号的频率为$f=25$kHz，红光的波长$\lambda\approx700$nm，因此可计算得到输出光功率
$$P_o=\eta f\dfrac{hc}{\lambda}=7.09\times10^{-16}\text{W}$$

使用光功率计可以测量得到发射一路激光时从Alice端出射的光功率约为$P_i=200\mu$W，因此可计算得到需要的衰减为
$$\alpha=10\lg\dfrac{P_i}{P_o}=114.5\text{dB}$$

实际上，由于发射光为具有间歇的脉冲，$P_i$会比200$\mu$W小，且输出端耦合也会造成一定的衰减，因此实际上需要的衰减量会比114.5dB小一些，可以根据单光子计数器测量的结果适当调整。

衰减片放置意味着光路系统已经完全完成，可以遮光进行实验测试。使用多层蒙布进行遮光，能够基本完全消除环境光对单光子计数的影响。值得注意的一点是由于偏振分束镜的不理想性，有一部分光会直接从输入端的偏振分束镜泄漏至输出端，对测量结果造成影响，因此需要在这一路径上单独遮光。

\subsubsection{同步方案与接收端解码}
% 区间处理，以及T2模式转csv再生成原始密钥的过程
实际通信时输入端与输出端是难以直接物理互连的。但是在我们的实验中，由于实验条件的限制，在发送光信号中添加标记信号或时钟信号再从输出端解码不太现实。因此做出妥协，将输入端的发送使能信号直接接至单光子计数器作为同步信号。

实验中测得暗计数率$R_d\approx600\cdot\text{s}^{-1}$，若不对光子计数数据进行筛选，则由暗计数造成的通信误码率为(估计$\nu\eta\approx0.1$)(估算时认为每一次暗计数都会造成一个误码)：
\[E=\frac{R_d}{R_b\mu\eta}\approx24\%\]

该结果无法被接受，考虑到每个比特实际上对应一个短脉冲，可以使用类似于“门控”的思路，只在激光器发光时进行单光子的探测。

我们实际上并无相应的门控元件，因此需要在数据后处理时对单光子事件进行筛选，筛选出激光器实际发光时接收到的单光子信号。

考虑到激光器的建立时间，单光子探测器的响应时间等因素，接收到发送使能信号与接收端上位机接收到脉冲对应的光子信号之间有一定时间差，且该时间差难以从理论上进行估计，为确定该时间差，发送3000kbit的信息，统计接收端的“单光子事件”与使能信号的接收时间差的分布，如图：
\begin{figure}[H]
    \centering
    \includegraphics[width=10cm]{graph5.png}
    \caption{单光子事件与使能信号的接收时间差分布图}
    %\label{fig:my_label}
\end{figure}

由图可见，光脉冲会在使能信号到达的$(6\sim11)\times10^6$ps后被接收。考虑到$7\times10^6$ps前接收到的脉冲光强较弱，数据后处理时只筛选出单周期内仅有一个单光子且落在使能信号到达后的的$(7\sim11)\times10^6$ ps内的事件作为“有效事件”。此时暗计数带来的误码率约为(式中$\tau=4×10^6$ ps，为“门”的宽度)：
\[E=\frac{{\tau R}_d}{\mu\eta}\approx2.4\%\]

这一误码率相对较低，可以被接受。

实际操作中，我们使用单光子探测器的 T2 模式，将每个单光子事件记录为「绝对时间戳 + 通道号」二元组，同步信号则作为一路特殊的「通道 0 事件」与光子事件混合记录在同一条时间轴上。实验结束后从单光子计数器导出原始 \texttt{.txt} 文件，通过 Python 脚本进行如下处理：

\begin{enumerate}
    \item 按时间顺序遍历文件，以同步通道的事件为边界将整条时间轴划分为若干「周期窗口」，每个窗口对应 Alice 发送的一个字节；
    \item 在每个窗口内统计落在 $(7\sim11)\times10^6$ ps 时间区间（即门控窗口）内的光子事件；若窗口内此区间事件数不为 1，则丢弃该窗口（多光子事件或无事件无法成码）；
    \item 对每个保留下来的事件，根据通道号查表得到 Bob 端的偏振索引；将偏振索引对 2 取模即得到比特值（0/2 → 0, 1/3 → 1），同时偏振索引除以 2 即得到所用的接收基；
    \item 将上述「序号、Alice 偏振索引、Bob 偏振索引、是否成码」四列写入 CSV 文件，作为后续基比对、QBER 估算与 LDPC 纠错的统一输入。
\end{enumerate}

这一处理流程同时完成了「同步对齐 + 门控筛选 + 通道映射」，是原始单光子事件转换为可用 raw key 的关键环节。

\subsubsection{数据后处理}

% 这一段个人认为可以侧重于代码实现
\textbf{基比对}

接收端解码后，发送端与接收端使用经典信道交换发送/接收时选择的测量基。扔掉基不匹配的测量结果，保留基一致的测量结果用于成码。实际实验中未模拟双方通过经典信道交换测量基的过程，而是直接将数据汇总到一台电脑上进行比对。

\textbf{真实误比特率计算}

基比对后，每个匹配的单光子事件可成码 1  比特。为衡量到此步骤为止系统的性能，直接在同一台电脑上对双方的单光子事件进行成码+直接比对\footnote{本步骤仅为了对系统的指标进行测量。实际 QKD 中发送端与接收端的单光子测量结果和成码后的数据不能通过经典信道传输，因此无法同时获得两端的原始数据，不能进行直接比对。}，得到真实的 QBER。结果如下表所示。

\begin{table}[!ht]
    \centering
    \caption{系统QBER测试结果数据表}
    \begin{tabular}{|l|l|l|l|l|}
    \hline
        实验编号 & 发送数据/kbit & 成码数据/bit & 正确成码/bit & QBER \\ \hline
        1 & 750 & 32892 & 29604 & 0.099991 \\ \hline
        2 & 1500 & 50519 & 45412 & 0.101091 \\ \hline
        3 & 3000 & 129827 & 116206 & 0.104917 \\ \hline
    \end{tabular}
\end{table}

可见，实验中QBER<0.11，可以生成安全密钥。

\textbf{纠错}

在接收端成码后，使用低密度奇偶校验码（LDPC）配合多轮披露进行纠错。具体流程为：将原始密钥按固定长度（实验中取 $N=2520$ 比特）分组，对每个分组按估计的 QBER 选择合适码率的校验矩阵，发送端通过经典信道发送校验子，接收端用置信传播算法迭代求解错误模式；若单次迭代未能成功，则公开少量「最不确定」的比特真值作为辅助信息，再次迭代，直至所有分组均成功对齐。整套算法在一台电脑上用 Python 模拟（依赖 \texttt{numpy}、\texttt{scipy}、\texttt{sympy}，代码见 \texttt{ecc/ecc\_simplified.py}），未模拟通过经典网络传输校验子和披露比特的网络传输部分。

在发送数据量为 $750$ kbit 的实验中，所有分组均成功收敛，双方持有的密钥完全一致（即纠错后 QBER 严格为 $0$）；纠错过程的平均效率为 $0.7226$\footnote{此处效率为 $0\sim1$ 之间的归一化值，$1$ 为理想情况（即实际泄漏量恰好等于 Shannon 理论下界，纠错最高效）；值越小表示实际泄漏的信息越多于理论必需量。}。

\textbf{安全性分析+私密放大}

在QKD实验中，原则上需要通过安全性分析计算全部流程中泄露信息量的上界，并通过私密放大将纠错后密钥进行压缩得到安全密钥，以使得窃听者无法从泄露的信息中得到安全密钥的任何信息，实现信息论意义上的安全。但由于由于安全性分析所需数据量较大，因此没有在实验中包含此部分。
\section{实验结果}

在本次实验中，我们成功搭建了偏振编码BB84协议量子密钥分发实验系统，并进行了量子密钥的传输与后处理纠错。我们最终进行了 750 kbit 数据的传输，接收端基比对后得到的 raw key 比特数约为 33 kbit\footnote{实验中精确值为 32892 bit。}，原始 QBER 约为 $10\%$\footnote{实验中精确值为 $9.9991\%$。}。所有分组的 LDPC 纠错均成功收敛，纠错后双方持有的密钥完全一致（QBER $=0$），实际比特率约为 1 kbit$\cdot$s$^{-1}$。

\section{分析与讨论}

\subsection{暗计数造成的误码}

暗计数造成的误码率已经在“接收端解码”这一节进行了估算，此处不对计算过程进行重复赘述。最终估算出：不考虑其他因素的影响，那么暗计数自身造成的误码率应在$2.4\%$左右。

\subsection{偏振分束镜不理想性造成的误码}

实验中，理论解码概率（发送的不同偏振态光子被不同探测器接收到的概率）如下表所示：

\begin{table}[!ht]
    \centering
    \caption{理论解码概率}
    \begin{tabular}{|l|l|l|l|l|}
    \hline
        \diagbox{接收}{发送} & 00 & 01 & 10 & 11 \\ \hline
        00 & 0.5 & 0 & 0.25 & 0.25 \\ \hline
        01 & 0 & 0.5 & 0.25 & 0.25 \\ \hline
        10 & 0.25 & 0.25 & 0.5 & 0 \\ \hline
        11 & 0.25 & 0.25 & 0 & 0.5 \\ \hline
    \end{tabular}
\end{table}

而在实际实验中，由于光路的不理想性(实验中使用的偏振分束镜并不完全理想，且光路中偏振片调节的精度有限)，即使接收端选择了正确的基，发射出的光子也有可能被错误的探测器接收。下表是（未调节各路光强，未安放衰减片时）单独点亮各个激光器时，各个接收器的光强数据表格(表中单位均为$\mu$W)：

\begin{table}[!ht]
    \centering
    \caption{单独点亮各激光器时接收光强数据表}
    \begin{tabular}{|l|l|l|l|l|}
    \hline
        \diagbox{接收}{发射} & 00 & 01 & 10 & 11 \\ \hline
        00 & 216 & 21 & 276 & 107 \\ \hline
        01 & 20 & 227 & 277 & 91 \\ \hline
        10 & 113 & 119 & 434 & 31 \\ \hline
        11 & 114 & 131 & 71 & 163 \\ \hline
    \end{tabular}
\end{table}

据此可以计算出实际解码概率：

\begin{table}[!ht]
    \centering
    \caption{实际解码概率}
    \begin{tabular}{|l|l|l|l|l|}
    \hline
        \diagbox{接收}{发送} & 00 & 01 & 10 & 11 \\ \hline
        00 & 0.43 & 0.04 & 0.23 & 0.30 \\ \hline
        01 & 0.04 & 0.39 & 0.19 & 0.37 \\ \hline
        10 & 0.26 & 0.23 & 0.43 & 0.08 \\ \hline
        11 & 0.12 & 0.26 & 0.06 & 0.57 \\ \hline
    \end{tabular}
\end{table}

可见，在进行基筛选后，解码正确与解码错误的光子数之比约为10:1.

需要注意的是，由于11通道探测器的耦合效率强于10通道，使得发送10时的误码率高于发送11时的误码率。同时，由于接收基不正确的结果会被直接舍弃，因此部分信道不平衡(如发送00时10和11的信道不平衡)对结果其实没有影响。

综合上述讨论，实验中纠错前的QBER应约为
\[1-(1-10\%)(1-2.4\%)=12.16\%\]

与实测QBER相近。实测QBER实际上小于$12.16\%$，这可能是因为在计算暗计数带来的误码率时对$\mu\eta$这个参数的估计偏小。

\section{体会建议}

我们的整个实验过程并不十分顺利。由于最开始的实验分工有些混乱，加上常常有“偏振片或激光器被不小心碰歪或转动”之类的意外发生，导致我们在最开始的光路调节环节中干了不少冤枉活，浪费了一些时间。后续我们吸取了教训，对已经调好的偏振片与分束镜的角度示数进行了记录，避免了一些重复劳动。但是，后续的等高共轴调节，以及输出端耦合的调节也十分痛苦，耗费了我们不少时间。同时，我们在驱动电路的搭建上也走了弯路。最初，我们打算使用光耦，但是我们发现其工作特性与我们的实验不匹配，最后，我们在老师的建议下选择了使用运算放大器来搭建驱动电路。除此之外，在对接收到的单光子数据进行后处理时，我们起初得到的结果与“瞎猜”没有明显的区别，这个问题困扰了我们一段时间，最后发现其来自于单光子信号与同步信号之间的时间差，我们对这一时间差进行了测量，并采取了“加窗”的方式降低了暗计数的影响。

但是，在这一实验中，这些不顺利也带给我们许多的收获。我们学习了PBS等器件的特性，熟练掌握了光路的调节方式，积累了Verilog相关代码的编写以及FPGA板子调试的经验，学习了驱动电路的搭建，等等。这些经验对于我们都十分宝贵。

我们对于课程也有一些建议。首先，我们希望课程组能实验中的各种仪器，如激光器等，提供datasheet，或者在明显处标记其型号。此外，我们也希望实验室能够标注一下支杆等各种实验器材的存放位置。

\appendix
\section{小组分工}

董皓彧：主要负责输入端FPGA控制与随机生成发射信号的代码编写，以及输入端驱动电路的搭建

刘佳起：主要负责光路搭建与联调，以及输出端信息解码的代码编写

周王昊：主要负责光路搭建与联调，以及答辩PPT与实验报告的框架搭建

\end{document}
